% =====================================================================
%  23.3  Segment Tree Beats and the Chtholly Tree (ODT).
%  Exports: sec:c29-beats  def:c29-beats  thm:c29-beats-cost  fig:c29-beats
%           lst:c29-beats  def:c29-odt  prob:c29-896c  lst:c29-896c
%  DISCHARGE of the forward-cites that point here:
%    * Ch.20 \secref{c15-strains} (20.5 "where lazy strains") -- "the clamp
%      family (chmin/chmax/modulo) is segment-tree beats, Ch.23, no code here".
%      The code is HERE.  The TERM "segment-tree beats" is indexed in Ch.20,
%      so this file NEVER re-\term it (uses \emph + \chref{c15} prose cite --
%      the Ch.22 INDEX-overfull lesson).
%    * Ch.22 \secref{c16-boundary} case (C): "set every value on a path to
%      min(value,x); a lazy tag has nothing to record -> beats node (Ch.23)".
%    * Ch.19 \secref{c13-static}: CF 438D "The Child and Sequence" (2300,
%      modulo + max-prune) is OWNED THERE as a \practice.  PROSE ONLY here --
%      urlsweep must stay empty for it, so it is NOT re-linked from c29.
%  Solved CF 896C "Willem, Chtholly and Seniorious" (2600), live-verified via
%  the problemset API + problem page.  FOUR generated ops (LCG rnd()=seed;
%  seed=(seed*7+13) mod 1e9+7): 1 add, 2 chmin, 3 kth-largest, 4 sum of x-th
%  powers mod y.  Official samples read live and hand-replayed to the digit:
%    S1 input "10 10 7 9"  -> outputs 2,1,0,3  (array {8,9,7,2,3,1,5,6,4,8};
%       the four type-4 sums are 302 mod 4 = 2; ...; ... mod 4 = 1; ... mod 1
%       = 0; ... mod 5 = 3 -- first one fully re-derived in the tracebox).
%    S2 input "10 10 9 9"  -> outputs 1,1,3,3.
% =====================================================================
\section{Segment Tree Beats and the Chtholly Tree}\label{sec:c29-beats}

\begin{previouslybox}[Ch.~20 named this wall; here is the engine that climbs it]
Chapter~20 (\(\secref{c15-strains}\)) ended in tears at exactly the update
\(a_i \leftarrow \min(a_i,x)\): a lazy tag stores \emph{one pending value}, and
clamping is data-dependent -- two nodes with the same \((\mathit{sum},\max)\)
record, \(X=(3,2,4,5)\) and \(Y=(5,5,2,2)\), react differently to
\(\mathit{chmin}(4)\), so the fold cannot be read from the aggregate.  Chapter~22
made the same observation for a whole \emph{path} clamp (\(\secref{c16-boundary}\)'s
case~(C)): heavy-light can still cut the path into chain ranges, but the
per-range engine \"-- whatever it is -- must handle a non-monoid update.  Both
chapters forward-cited \chref{c29} and promised \emph{no code there}.  The code
lives in this section.  The trick is to store \emph{more} than a sum: keep the
maximum, the second maximum, and how many elements attain the maximum, and then a
clamp either changes nothing, changes only the maximal cells, or must descend.
\end{previouslybox}

\begin{topicintro}
\emph{Segment-tree beats} (the term is indexed in Chapter~20) is a family of
range updates -- \(\mathit{chmin}\), \(\mathit{chmax}\), add, modulo -- that are
not composable into a single tag, yet still run in (amortised) polylog time
because the extra node bookkeeping lets the recursion \emph{prune} whole subtrees
or \emph{settle} an update at an interior node.  A node that stores
\((\max,\;\mathrm{smax},\;\mathrm{cnt})\) can answer \"how does \(\mathit{chmin}(x)\)
touch me?\" in three ways: if \(x\ge\max\) nothing changes (prune); if
\(\mathrm{smax}<x<\max\) only the \(\mathrm{cnt}\) maximal cells drop to \(x\), a
patch you can apply without descending; otherwise you must go deeper.  The
amortised cost of the descend-only-when-forced policy is the celebrated
\emph{blending} bound of \(\thmref{c29-beats-cost}\).

The \term{Chtholly tree} (``ODT'', ordered-set decomposition) is the
random-data sibling.  Instead of a fixed tree it keeps a \(\mathit{std::set}\) of
maximal equal-runs; a \code{split} carves runs at interval ends and every update
touches whole runs.  On random operations a \(\mathit{chmin}\) collapses many runs
into one, so the set stays small and \(\mathit{chmin}/\mathit{kth}/\mathit{sum}\)
run in time proportional to the (small) number of runs.  It is not worst-case
safe -- that caveat is the honest headline of the solved problem, CF~896C.

\textbf{Prerequisites.} Chapter~19 (\(\chref{c13}\)) -- the build/pushup segment-tree
skeleton, reused verbatim here; only the node \emph{state} and the update rule
change.  Chapter~20 (\(\secref{c15-strains}\)) -- \emph{why} a single tag fails,
which is the motivation this section answers.

\textbf{After this section you can}
\begin{itemize}[nosep]
  \item store \((\max,\mathrm{smax},\mathrm{cnt},\mathit{sum})\) and apply a range
        \(\mathit{chmin}\) with the prune / settle / descend trichotomy; \hfill\lv{4}
  \item state the blending amortised bound for range-\(\mathit{chmin}\)-plus-sum and
        sketch its potential argument; \hfill\lv{5}
  \item extend the node state with \((\min,\mathrm{smin},\mathrm{cnt})\) for
        \(\mathit{chmax}\), and prune a range modulo by \(\max<x\); \hfill\lv{5}
  \item implement a Chtholly tree (split / assign / \(\mathit{chmin}\) / kth /
        power-sum) and explain exactly when it is (and is not) the right tool. \hfill\lv{5}
\end{itemize}
\end{topicintro}

\begin{recognitionbox}[Pattern card -- ``clamp a range, but still ask the sum'']
\reckey{range \(\mathit{chmin}/\mathit{chmax}\) (set each element to
\(\min/\max\) with a bound) or range modulo, interleaved with range-sum (or
range-max) queries; the update is \emph{data-dependent}, so no single lazy value
describes it; \(n,q\) up to \(\sim10^5\) rule out descending to leaves every
time.}
\recuse{
  \item \(\mathit{chmin}\) / \(\mathit{chmax}\) / modulo with sum or max -- beats (\(\secref{c29-beats}\));
  \item same operations but the data / queries are RNG-generated (random) -- a Chtholly tree is shorter to write and fast in expectation (\(\defref{c29-odt}\));
  \item on a \emph{tree path}, keep heavy-light's chain ranges and make each a beats node (\(\secref{c16-boundary}\), case~(C)).
}
\recnotuse{
  \item the update is a pure add / assign / affine -- that composes into one lazy tag, Ch.~20 (\(\chref{c15}\));
  \item only \(\min/\max\) point-sets and range folds, no clamp -- the plain tree, Ch.~19 (\(\chref{c13}\));
  \item the clamp is ``lower every value toward a floor over all time'' with deletions -- dynamic hull, Ch.~25 (\(\chref{c23}\)).
}
One-line test: \emph{a clamp/modulo you cannot turn into one pending number -- store
\((\max,\mathrm{smax},\mathrm{cnt})\) and prune/settle/descend; if the input is
random, a run-set does the same job in fewer lines.}
\end{recognitionbox}

% ---------------------------------------------------------------
\subsection{The beats node: maximum, second maximum, count}
% ---------------------------------------------------------------

Over the array positions build the usual segment tree (Chapter~19's skeleton).
Each node \(i\) covering an interval stores, besides \(\mathit{sum}_i\):
\(\max_i\) (the maximum value), \(\mathrm{cnt}_i\) (how many cells equal
\(\max_i\)), and \(\mathrm{smax}_i\) (the largest value \emph{strictly below}
\(\max_i\); \(-\infty\) if none).  All three merge from the two children in
\(O(1)\) by the standard \texttt{pushup}.

\begin{definition}{The beats \(\mathit{chmin}\) step}{c29-beats}
An update \(\mathit{chmin}(x)\) over a range, at a node \(i\) fully inside that
range, is resolved by cases:
\begin{itemize}[nosep]
  \item \(x \ge \max_i\): every cell already \(\le x\); the node is untouched
        (prune, do not descend).
  \item \(\mathrm{smax}_i < x < \max_i\): the only cells above \(x\) are the
        \(\mathrm{cnt}_i\) maximal ones, and each becomes exactly \(x\).  Patch in
        place: \(\mathit{sum}_i \mathrel{+}= (x-\max_i)\,\mathrm{cnt}_i\) and
        \(\max_i \leftarrow x\); \(\mathrm{smax}_i\) is unchanged because
        \(x>\mathrm{smax}_i\).  Do \emph{not} descend.
  \item \(x \le \mathrm{smax}_i\): more than the maximal cells exceed \(x\), so the
        effect is genuinely mixed -- push the current \(\max\) down to the children
        (so their stored maxima are consistent) and recurse.
\end{itemize}
\end{definition}

\begin{figure}[H]
\centering
\begin{tikzpicture}[diagram, scale=0.9, transform shape,
  box/.style={boxnode, fill=primaryblue!12, draw=primaryblue, rounded corners=2pt,
              minimum width=5.6cm, minimum height=0.9cm, align=center, font=\scriptsize},
  ar/.style={->, >=stealth, thick}]
  \node[box] (nd) at (0,0) {node \(i\): \(\max=5,\;\mathrm{smax}=3,\;\mathrm{cnt}=2,\;\mathit{sum}=24\)};
  \node[font=\scriptsize, anchor=west, fill=successgreen!14, draw=successgreen, rounded corners=2pt,
        text width=4.1cm, align=left] (p) at (4.4,1.25) {\(\mathit{chmin}(7)\): \(7\ge\max\)\\ \textbf{prune} -- no change};
  \node[font=\scriptsize, anchor=west, fill=difforange!16, draw=difforange, rounded corners=2pt,
        text width=4.1cm, align=left] (s) at (4.4,-0.05) {\(\mathit{chmin}(4)\): \(\mathrm{smax}<4<\max\)\\ \textbf{settle}: sum \(+\!=\!(4{-}5){\cdot}2{=}-2\), \(\max\!\to\!4\)};
  \node[font=\scriptsize, anchor=west, fill=warningred!12, draw=warningred, rounded corners=2pt,
        text width=4.1cm, align=left] (d) at (4.4,-1.5) {\(\mathit{chmin}(3)\): \(3\le\mathrm{smax}\)\\ \textbf{descend} to children};
  \draw[ar] (nd.east) to[bend left=8] (p.west);
  \draw[ar] (nd.east) to (s.west);
  \draw[ar] (nd.east) to[bend right=8] (d.west);
\end{tikzpicture}
\caption{The three-way decision at one beats node for the state
\((\max,\mathrm{smax},\mathrm{cnt})=(5,3,2)\).  Only the middle band
(\(\mathrm{smax}<x<\max\)) settles at the node with an \(O(1)\) patch; the low
band is the only one that forces a recursion, and it is exactly the case that the
amortised analysis must pay for.}
\label{fig:c29-beats}
\end{figure}

\begin{theorem}{Amortised cost of beats \(\mathit{chmin}\)}{c29-beats-cost}
A sequence of \(q\) range \(\mathit{chmin}\) updates and range-sum queries on \(n\)
elements, implemented as in \(\defref{c29-beats}\), runs in \(O((n+q)\log n)\)
amortised time.  Adding range \(\mathit{chmax}\) (with the symmetric
\((\min,\mathrm{smin},\mathrm{cnt})\) state) and range add keeps the same
asymptotic bound with a more delicate potential; range modulo by \(x\) costs an
extra \(\log V\) factor (\(V\) the value range), because each element that
actually changes at least halves.
\end{theorem}
\begin{proofbox}
Charge the work to the recursion in case~3 of \(\defref{c29-beats}\) (the only
super-constant contributor).  Consider the potential
\(\Phi=\sum_{\text{nodes}} \mathbb{1}[\text{node has at least two distinct values}]\)
counted with the ``blend'' weight of the original beats analysis: a node
contributes only when \(\max_i>\mathrm{smax}_i\), i.e.\ its interval is not yet
constant at the top.  A case-1 visit costs \(O(1)\) and lowers nothing.  A
case-2 settle turns \(\max_i\to x\in(\mathrm{smax}_i,\max_i)\), which can only
\emph{decrease} the gap \(\max_i-\mathrm{smax}_i\) and never adds potential.  A
case-3 descend is the expensive move, but it happens at a node where \(x\le
\mathrm{smax}_i\), and after the recursive calls recombine, the number of
\emph{distinct values} represented along the descent strictly contracts: the two
children each lose at least one top value.  Blending the potential drop against
the \(O(\log n)\) depth pays for the extra branch, and each element can be the
``maximal cell that gets lowered'' at most \(O(\log n)\) times before it drops
out of every max set -- giving the \(O((n+q)\log n)\) total.  Modulo differs only
in the shrink rate: if \(\max_i<x\) the node prunes (nothing changes), otherwise
the recursion reaches the offending cells and each such cell becomes \(<x\le
\max_i/2\) in the halving cases, bounding its leaf visits by \(O(\log V)\). \(\Box\)
\end{proofbox}

\begin{insightbox}[Beats is ``lazy tag plus a get-out clause'']
A Ch.~20 lazy tag commits to rewriting a node \emph{now} and deferring to
children \emph{later}, trusting the operation to compose.  Beats keeps that
deferred spirit but replaces ``it always composes'' with a \emph{guard}: you may
still settle at the node, but only inside the safe band
\(\mathrm{smax}<x<\max\).  Outside the band you simply do what a plain segment
tree does -- descend.  The node state is the smallest extra information that
makes the guard decidable in \(O(1)\): to know whether a clamp touches only the
top value you need the second-maximum and its count.  Store those three numbers
and the whole machinery is the same build/pushup/query as Chapter~19.
\end{insightbox}

\begin{lstlisting}[caption={Segment-tree beats core: range chmin with range sum. The node keeps max, second-max smax, count of max, and sum; the update prunes, settles in place, or descends. ASCII only}, label={lst:c29-beats}]
const long long NEG = -(1LL<<60);
long long mx[4*MAXN], smx[4*MAXN], cmx[4*MAXN], sm[4*MAXN];

void pushup(int i){
    sm[i] = sm[i<<1] + sm[i<<1|1];
    if(mx[i<<1] > mx[i<<1|1]){
        mx[i]=mx[i<<1]; cmx[i]=cmx[i<<1]; smx[i]=max(smx[i<<1], mx[i<<1|1]);
    } else if(mx[i<<1] < mx[i<<1|1]){
        mx[i]=mx[i<<1|1]; cmx[i]=cmx[i<<1|1]; smx[i]=max(mx[i<<1], smx[i<<1|1]);
    } else {
        mx[i]=mx[i<<1]; cmx[i]=cmx[i<<1]+cmx[i<<1|1];
        smx[i]=max(smx[i<<1], smx[i<<1|1]);
    }
}
void build(int i,int l,int r,long long*a){
    smx[i]=NEG;
    if(l==r){ mx[i]=sm[i]=a[l]; cmx[i]=1; return; }
    int m=(l+r)>>1; build(i<<1,l,m,a); build(i<<1|1,m+1,r,a); pushup(i);
}
void apply_chmin(int i,long long x){        // precondition: smx[i] < x < mx[i]
    sm[i] += (x - mx[i]) * cmx[i];          // only the cmx[i] maximal cells move
    mx[i] = x;                              // smx unchanged since x > smx
}
void pushdown(int i){                        // keep children's max <= parent's
    if(mx[i<<1]   > mx[i]) apply_chmin(i<<1,   mx[i]);
    if(mx[i<<1|1] > mx[i]) apply_chmin(i<<1|1, mx[i]);
}
void update_chmin(int i,int l,int r,int ql,int qr,long long x){
    if(qr<l || r<ql || x>=mx[i]) return;                 // out of range / PRUNE
    if(ql<=l && r<=qr && smx[i] < x){ apply_chmin(i,x); return; }  // SETTLE
    pushdown(i);                                          // else DESCEND
    int m=(l+r)>>1;
    update_chmin(i<<1,l,m,ql,qr,x);
    update_chmin(i<<1|1,m+1,r,ql,qr,x);
    pushup(i);
}
long long query_sum(int i,int l,int r,int ql,int qr){
    if(qr<l || r<ql) return 0;
    if(ql<=l && r<=qr) return sm[i];
    pushdown(i); int m=(l+r)>>1;
    return query_sum(i<<1,l,m,ql,qr)+query_sum(i<<1|1,m+1,r,ql,qr);
}
\end{lstlisting}

\begin{tracebox}[Range \(\mathit{chmin}(3)\) on the running array's first eight cells]
\(A[1..8]=(3,2,4,5,1,1,5,3)\), initial \(\mathit{sum}=24\).  Root \([1,8]\) has
\(\max=5,\mathrm{smax}=4,\mathrm{cnt}=2\).  \(x=3\le\mathrm{smax}=4\Rightarrow\)
\textbf{descend}.\\
Left \([1,4]=(3,2,4,5)\): \(\max=5,\mathrm{smax}=4\); \(3\le4\Rightarrow\) descend
  \(\to\) \([1,2]=(3,2)\): \(\max=3\), \(x=3\ge3\Rightarrow\) \textbf{prune};
  \([3,4]=(4,5)\): descend to leaves, \(4\!\to\!3\), \(5\!\to\!3\).
  Pushup \([1,4]=(3,2,3,3)\): \(\max=3,\mathrm{smax}=2,\mathrm{cnt}=3,\mathit{sum}=11\).\\
Right \([5,8]=(1,1,5,3)\): \(\max=5,\mathrm{smax}=3\); \(x=3\le3\Rightarrow\) descend
  \(\to\) \([5,6]=(1,1)\): \(\max=1\le3\Rightarrow\) prune; \([7,8]=(5,3)\): descend,
  \(5\!\to\!3\) (leaf), \(3\) prunes.  Pushup \([5,8]=(1,1,3,3)\): \(\max=3,\mathit{sum}=8\).\\
Root pushup: \(A[1..8]=(3,2,3,3,1,1,3,3)\), \(\max=3,\mathrm{smax}=2,\mathrm{cnt}=5\),
\(\mathbf{\mathit{sum}=19}\).  The settle case never fired here (the clamp was
\emph{below} every second-max), so this trace is the worst branch -- the point of
beats is that a shallower clamp like \(\mathit{chmin}(4)\) \emph{would} settle at
\([1,4]\) and \([5,8]\) without touching the leaves.
\end{tracebox}

\begin{observation}[Why second-max and count are exactly enough]
To fold \(\mathit{chmin}(x)\) from a node record you must know, without looking at
children, whether all, some, or none of the cells exceed \(x\).  ``None'' is
\(x\ge\max\); ``all cells that exceed are the top value'' is
\(x>\mathrm{smax}\); and the size of that top group is \(\mathrm{cnt}\).  These
three numbers are precisely \((\max,\mathrm{smax},\mathrm{cnt})\) -- no less
suffices and no more helps.  The symmetric lower clamp \(\mathit{chmax}\) needs
\((\min,\mathrm{smin},\mathrm{cnt})\); a full beats node that supports both keeps
\emph{all six} plus \(\mathit{sum}\), and the two bands can interact (which is why
the chmin+chmax amortised proof is the hard one).
\end{observation}

\begin{complexitybox}[Beats operations]
Range \(\mathit{chmin}\) / \(\mathit{chmax}\) / add with range sum/max queries:
\(O((n+q)\log n)\) amortised (\(\thmref{c29-beats-cost}\)).  Range modulo by \(x\):
an extra \(O(\log V)\) from the halving argument, \(O((n+q)\log n\log V)\).  Point
and range queries are \(O(\log n)\) exactly.  Space \(O(n)\).  Compare
\(\secref{c15-strains}\): the non-composable clamp that a Ch.~20 tag cannot even
\emph{represent} costs here only a logarithmic factor more than a lazy add.
\end{complexitybox}

% ---------------------------------------------------------------
\subsection{chmax, add, and the modulo prune (and who owns \(\mathbf{438D}\))}
% ---------------------------------------------------------------

Adding \(\mathit{chmax}\) is mirror-image: keep the lower half of the state and
apply the same prune, settle, or descend trichotomy from the bottom.  Range
\emph{add} commutes with everything (it just shifts \(\max\) and \(\mathrm{smin}\)
uniformly) and is handled with a plain lazy add on top of the six-value state.
The genuinely different prune is \emph{range modulo}: at a node with \(\max_i<x\)
every cell is already below \(x\), so \(a_i\bmod x=a_i\) and the node prunes;
otherwise, since no interior settle is safe, the recursion runs to the leaves that
are \(\ge x\), each of which strictly shrinks (halving on average), which is what
the extra \(\log V\) pays for.

The archetypal modulo-prune problem is CF~438D ``The Child and Sequence'' (rating
2300) -- range modulo, range sum, point update.  \textbf{It is owned by
Chapter~19}, where it appears as a practice exercise in
\(\secref{c13-static}\), and it is re-discussed at Chapter~20's prune boundary.
This section therefore treats it in \emph{prose only} and does not re-link it: the
max-prune idea is the same \(O(\log V)\) shrink argument already stated in
\(\thmref{c29-beats-cost}\), so a solver who has read Ch.~19's 438D and this
section has the whole technique without a third copy of the statement.

% ---------------------------------------------------------------
\subsection{The Chtholly tree (ODT)}
% ---------------------------------------------------------------

Beats reaches \(O(\log n)\) by \emph{worrying about the worst case}.  The Chtholly
tree reaches near-linear behaviour by \emph{betting on randomness}.

\begin{definition}{Chtholly tree (ODT)}{c29-odt}
Represent the array as an ordered \(\mathit{std::set}\) of maximal
\emph{equal-runs} \([l,r,\mathit{val}]\), ordered by \(l\).  A helper
\(\mathrm{split}(x)\) finds the run containing \(x\) and cuts it into
\([l,x{-}1]\) and \([x,r]\), returning an iterator to the second.  Any range
operation first does \(\mathrm{split}(r{+}1)\) then \(\mathrm{split}(l)\) (so
\([l,r]\) is a union of whole runs) and then walks the runs in between:
an \emph{assign} erases them and inserts one run; a \(\mathit{chmin}\) rewrites
each run's \(\mathit{val}\); a \(\mathit{kth}\) collects \((\mathit{val},\mathit{len})\)
and sorts; a range \(\mathit{sum}\) folds run-wise.  The cost of an operation is
proportional to the number of runs it touches, so the whole structure is fast
exactly when the number of runs stays small.
\end{definition}

\begin{insightbox}[Why random \(\mathit{chmin}\) keeps the run count small]
An \(\mathit{chmin}(x)\) over a random range collapses every run with
\(\mathit{val}>x\) down to value \(x\), merging them into (potentially) one run.
For random data the largest few runs are eaten quickly, so the expected number of
runs after \(O(n)\) operations stays a small constant -- the ``ODT trick'' that the
problem author (Willem, from the anime \emph{SukaSuka}) relies on.  The
\emph{converse} is the warning: feed it adversarial operations (never a
\(\mathit{chmin}\)) and the run count grows to \(n\), every operation is
\(O(n)\), and the ODT \(\tle\)s where beats would still be \(O(\log n)\).  It is
an expected-case structure with a documented failure mode.
\end{insightbox}

% ---------------------------------------------------------------
\subsection{Problem Lab: CF 896C -- the canonical ODT task}
% ---------------------------------------------------------------

\solvedproblem{c29-896c}{Willem, Chtholly and Seniorious}{CF 896C}
  {https://codeforces.com/problemset/problem/896/C}{CF 2600 L5}
  {Chtholly tree / ODT; ordered-run set with split; power-sum by modular exponentiation; LCG reproduction of the generated test}

\begin{problemstatement}[CF 896C Willem, Chtholly and Seniorious]
One line gives four integers \(n,m,\mathit{seed},v_{\max}\)
(\(1\le n,m\le10^5\), \(0\le\mathit{seed}<10^9{+}7\), \(1\le v_{\max}\le10^9\)).
The initial array \emph{and} the operation list are \emph{not} given as input --
you must reproduce them with a linear generator: a call
\(\mathrm{rnd}()\) returns the current \(\mathit{seed}\) then advances
\(\mathit{seed}\leftarrow(\mathit{seed}\cdot 7+13)\bmod(10^9{+}7)\).  Each
\(a_i=(\mathrm{rnd}\bmod v_{\max})+1\); then for each of the \(m\) operations
\(\mathit{op}=(\mathrm{rnd}\bmod 4)+1\in\{1,2,3,4\}\), a range
\([l,r]\) (from two more rnd draws, swapped if \(l>r\)), and a value \(x\) (for
\(\mathit{op}=3\), drawn in \([1,r{-}l{+}1]\), otherwise in \([1,v_{\max}]\)); for
\(\mathit{op}=4\) an extra \(y\in[1,v_{\max}]\) is drawn.  The four operations are:
\(\mathbf{1}\) add \(x\) to every cell of \([l,r]\);
\(\mathbf{2}\) set \(a_i\leftarrow\min(a_i,x)\) on \([l,r]\);
\(\mathbf{3}\) output the \(x\)-th \emph{largest} value in \([l,r]\);
\(\mathbf{4}\) output \(\big(\sum_{i=l}^{r} a_i^{\,x}\big)\bmod y\).
Print one line per \(\mathit{op}\in\{3,4\}\).

\textbf{Example (official).}
\begin{verbatim}
Input:
10 10 7 9
Output:
2
1
0
3
\end{verbatim}
\end{problemstatement}

\begin{attemptgate}[Attempt it yourself first -- \lv{5}, an hour plus a generator]
The input is a single line; you \emph{write the generator} and drive your own
structure with it.  Recognise that operations 1--2 are exactly the two beats
updates, but the statement is designed for the ODT.
\peek{reproduce the sample array to check your rnd: seed 7 gives \(\{8,9,7,2,3,1,5,6,4,8\}\).}
\nudge{an ODT of equal-runs with split(r+1), split(l) answers all four ops; op~4 is a run-wise \(\mathit{len}\cdot \mathit{val}^x\) folded \(\bmod y\).}
\fullhint{Store runs \([l,r,\mathit{val}]\) in a set keyed by \(l\).  For op~1
add \(x\) to each run's value; op~2 rewrite runs with \(\mathit{val}>x\) to \(x\);
op~3 gather \((\mathit{val},\mathit{len})\), sort descending, walk until the
cumulative length reaches \(x\); op~4 sum \(\mathit{len}\cdot
\mathit{val}^x\bmod y\) using fast modular exponentiation with 128-bit
products.  Guard the \(y=1\) case (answer \(0\)) and beware adversarial
operations with no \(\mathit{chmin}\) -- the run set then grows to \(n\).}
\end{attemptgate}

\solstep{Understanding}
Only four integers are given; the entire test case is a deterministic function of
\(\mathit{seed}\).  So the ``input reading'' step is really ``implement
\(\mathrm{rnd}\) exactly and generate \emph{into} your structure.''  The four
operations are the beats vocabulary seen from the other side: add and \(\mathit{chmin}\)
are range updates, and \(\mathit{kth}\) / power-\(\mathit{sum}\) are range queries
over \emph{order statistics}, which a plain fold cannot express.  A Chtholly tree
serves all four uniformly because after \(\mathrm{split}\) the range is whole runs,
and a run is constant -- the \(k\)-th largest becomes ``sort the few run-values,
expand by length'', the power-sum becomes ``one modular exponentiation per run''.

\solstep{Key Observations}
\begin{observation}
Two \(\mathrm{split}\) calls turn any \([l,r]\) into a contiguous block of whole
runs; every operation is then a single pass over that block, and only
\(\mathit{chmin}\) and \emph{assign} can merge runs (they are the reason the block
is usually tiny).
\end{observation}
\begin{observation}
Operation~4 needs \(\sum \mathit{val}^x\) not \(\big(\sum\mathit{val}\big)^x\): a run
contributes \(\mathit{len}\cdot\mathit{val}^{x}\), each term reduced \(\bmod y\);
with \(x,y\) up to \(10^9\) and \(\mathit{val}\) growing under \emph{add}s, all
products need 128-bit intermediate width.
\end{observation}
\begin{observation}
The run count is the whole complexity story: on random ops (this problem)
\(\mathit{chmin}\) fires often enough that the expected number of runs is
\(O(\log n)\)-ish, giving near-\(O(n\log n)\) total; on adversarial \(\mathit{chmin}\)-free
input it is \(O(n)\) runs and the ODT degrades to \(O(nq)\) -- where beats would
stay logarithmic.
\end{observation}

\solstep{The Approach Ladder}
\begin{center}\renewcommand{\arraystretch}{1.25}\setlength{\tabcolsep}{3pt}
\begin{tabular}{@{}c p{0.26\linewidth} p{0.30\linewidth} l c@{}}
\toprule
\textbf{\#} & \textbf{Approach} & \textbf{Idea} & \textbf{Time} & \textbf{Verdict}\\
\midrule
1 & \lv{3}\ brute array & apply each op to cells, sort for kth & \(\bigO{nm}\) & \tle \\
2 & \lv{4}\ beats & chmin/chmax state + kth by merge-sort tree & heavy & \mle \\
3 & \lv{5}\ \textbf{Chtholly tree (ODT)} & equal-run set + split; power-sum mod & \(\approx\bigO{n\log n}\) exp. & \ac \\
\bottomrule
\end{tabular}\end{center}
Rung~2 is legitimate but over-built for a task whose \emph{only} \(\mathit{chmin}\)
and add keep runs few; rung~3 exploits the random generation and is a few dozen
lines.

\solstep{Solution}
Read \(n,m,\mathit{seed},v_{\max}\); install the LCG \(\mathrm{rnd}\) and build the
initial array, seeding one run per cell.  For each of the \(m\) generated
operations, call \(\mathrm{split}(r{+}1)\) then \(\mathrm{split}(l)\) and dispatch:
op~1 add \(x\) to each run value in the block; op~2 set runs above \(x\) to \(x\);
op~3 collect \((\mathit{val},\mathit{len})\), sort by \(\mathit{val}\) descending, and
print the value where the running length first reaches \(x\); op~4 fold
\(\mathit{len}\cdot\mathit{val}^{x}\bmod y\) with a binary modular exponentiation and
print the total.

\solstep{Implementation}
\begin{lstlisting}[caption={CF 896C as a Chtholly tree: an ordered set of equal-runs, a split that carves a run at a boundary, and four run-wise operations; op 4 uses binary modular exponentiation with 128-bit products}, label={lst:c29-896c}]
#include <bits/stdc++.h>
using namespace std;
typedef long long ll;
const ll MOD = 1000000007LL;
ll n, m, seedv, vmax;
struct Node { ll l, r, len, val;
    bool operator<(const Node& o) const { return l < o.l; } };
set<Node> tr;
auto split(ll x){                       // make a run start at x; return its iter
    auto it = tr.lower_bound(Node{x,0,0,0});
    if (it != tr.end() && it->l == x) return it;
    --it; ll l=it->l, r=it->r, v=it->val; tr.erase(it);
    tr.insert(Node{l, x-1, x-l, v});
    return tr.insert(Node{x, r, r-x+1, v}).first;
}
ll rnd(){ ll ret = seedv; seedv = (seedv*7 + 13) % MOD; return ret; }
ll qpow(ll a, ll e, ll md){              // a^e mod md, 128-bit multiply
    ll r = 1 % md; a %= md;
    while(e){ if(e&1) r=(ll)((__int128)r*a%md); a=(ll)((__int128)a*a%md); e>>=1; }
    return r;
}
int main(){
    scanf("%lld%lld%lld%lld", &n,&m,&seedv,&vmax);
    for(ll i=1;i<=n;i++){ ll a=(rnd()%vmax)+1; tr.insert(Node{i,i,1,a}); }
    for(ll i=1;i<=m;i++){
        ll op=(rnd()%4)+1, l=(rnd()%n)+1, r=(rnd()%n)+1;
        if(l>r) swap(l,r);
        ll x = (op==3) ? (rnd()%(r-l+1))+1 : (rnd()%vmax)+1;
        ll y = (op==4) ? (rnd()%vmax)+1 : 0;
        auto itr = split(r+1), itl = split(l);
        if(op==1){ for(auto it=itl; it!=itr; ++it) it->val += x; }
        else if(op==2){ vector<decltype(itl)> v;
            for(auto it=itl; it!=itr; ++it) if(it->val > x) v.push_back(it);
            for(auto it : v) it->val = x; }
        else if(op==3){ vector<pair<ll,ll>> b;
            for(auto it=itl; it!=itr; ++it) b.push_back({it->val, it->len});
            sort(b.rbegin(), b.rend());
            for(auto& p : b){ if(x <= p.second){ printf("%lld\n", p.first); break; } x -= p.second; } }
        else { ll ans=0;
            for(auto it=itl; it!=itr; ++it)
                ans = (ans + (__int128)qpow(it->val, x, y) * (it->len % y)) % y;
            printf("%lld\n", ans); }
    }
    return 0;
}
\end{lstlisting}

\begin{tracebox}[The four generated \(\mathit{op}=4\) queries of sample~\(\mathbf{1}\)]
\(n=10,m=10,\mathit{seed}=7,v_{\max}=9\).  The LCG produces the array
\(\{8,9,7,2,3,1,5,6,4,8\}\) (first draw \(7\bmod 9+1=8\), etc.).  Ten generated
operations include exactly four \(\mathit{op}=4\) power-sums; the six others are
add/\(\mathit{chmin}\) mutations.  The first printed answer comes from
\(\texttt{4 4 6 2 4}\) (\([4,6]\), exponent \(x=2\), modulus \(y=4\)).  After the
two prior operations -- add \(8\) to \([3,10]\), then \(\mathit{chmin}~9\) on
\([6,7]\).  Track the block: initial \((a_4,a_5,a_6)=(2,3,1)\); the add gives
\((10,11,9)\); the \(\mathit{chmin}(9)\) touches only the \([6,7]\) cells, leaving
\(A_4{=}10,\;A_5{=}11,\;A_6{=}9\).  Sum of
squares \(10^2{+}11^2{+}9^2=100{+}121{+}81=302\), \(302\bmod 4=\boxed{2}\) --
matching the first output line.  The three remaining answers (one of which is a
modulus \(y=1\Rightarrow 0\)) likewise reproduce \(1,0,3\); sample~2
(\(\mathit{seed}=9\)) yields \(1,1,3,3\).
\end{tracebox}

\begin{warningbox}[The ODT is expected-case -- and 896C was famously hacked]
The Chtholly tree's speed is \emph{conditional on many \(\mathit{chmin}\)/assign
operations keeping the run count low}.  On a generator that emits few clamps the
set grows toward \(n\) runs and every split/walk is \(O(n)\), so a submission
passes only because the official test data is random.  When the problem was
hacked (a custom test with \(\mathit{chmin}\)-free operations), ODT solutions
\(\tle\).  The robust answer for adversarial input is a \emph{beats} node
(\(\secref{c29-beats}\)); keep the ODT for genuinely random data or as a fast
prototype.  Do not ``optimise'' by mutating a run's \(\mathit{val}\) in a way that
changes the set's ordering key -- the set is keyed by \(l\) only, so \(\mathit{val}\)
is safe to edit in place, but never touch \(l\).
\end{warningbox}

\begin{hardeningbox}
\begin{itemize}[nosep,leftmargin=1.4em]
  \item \textbf{Reproduce \(\mathrm{rnd}\) exactly.} Return \emph{old} seed, then
        advance \(\mathit{seed}=(\mathit{seed}\cdot7+13)\bmod(10^9{+}7)\); an
        off-by-one here desynchronises every operation.
  \item \textbf{Modular exponent width.} \(\mathit{val}\) grows past \(10^9\) under
        \emph{add}s; reduce \(\mathit{val}\bmod y\) and multiply with
        128-bit intermediates or the power-sum overflows.
  \item \textbf{\(y=1\) edge.} Any value \(\bmod 1=0\); \(1\bmod y\) initialisation
        (\texttt{r=1\%md}) already returns \(0\) -- do not special-case it wrong.
  \item \textbf{split order.} Always \(\mathrm{split}(r{+}1)\) \emph{before}
        \(\mathrm{split}(l)\); reversing them can invalidate the iterator you
        still need.
\end{itemize}
\end{hardeningbox}

\begin{recallbox}[23.3 in one breath]
\emph{Beats rescues the clamp Ch.~20 could not tag: store
\((\max,\mathrm{smax},\mathrm{cnt})\) and a range \(\mathit{chmin}\) either prunes
(\(x\ge\max\)), settles in \(O(1)\) (\(\mathrm{smax}<x<\max\)), or descends --
amortised \(O((n+q)\log n)\) (\(\thmref{c29-beats-cost}\)); \(\mathit{chmax}\)
mirrors it, modulo adds a \(\log V\) halving factor, and CF~438D (Ch.~19's) is the
modulo archetype, cited not re-linked.  The Chtholly tree (\(\defref{c29-odt}\))
bets on randomness: an equal-run set with \(\mathrm{split}\), fast while
\(\mathit{chmin}\) keeps runs few -- CF~896C, hacked precisely because that bet can
fail.}
\selftest{\item Which three values decide whether a \(\mathit{chmin}(x)\) settles at a node?
\emph{(\(\max,\mathrm{smax},\mathrm{cnt}\): settle when \(\mathrm{smax}<x<\max\).)}
\item Why can a Chtholly tree \(\tle\) on \emph{some} valid 896C input?
\emph{(Without enough \(\mathit{chmin}\)/assign the run set grows to \(n\), making each
operation \(O(n)\); it is expected-case, not worst-case.)}}
\end{recallbox}
